\chapter{Introduction}
\label{ch:introduction}

\section{Background of the Study}
\label{sec:background}

Electronic waste (e-waste), commonly referred to as Waste Electrical and Electronic
Equipment (WEEE), is one of the fastest-growing solid waste streams in the world. The
rapid advancement of technology, increasing consumer demand for electronic devices,
shortened product life cycles, and continuous innovation have significantly increased
the production and disposal of electrical and electronic equipment. Devices such as
mobile phones, laptops, desktop computers, televisions, printers, refrigerators,
batteries, and other electronic appliances become obsolete within relatively short
periods, leading to substantial growth in electronic waste generation. As illustrated in
Figure~\ref{fig:ewaste-growth}, the volume of electronic waste generated worldwide
has increased steadily over the past decade, highlighting the growing need for
effective collection, recovery, and recycling systems.

\begin{figure}[H]
  \centering
  \imageplaceholder{Growth of Electronic Waste Management chart (adapted from the Global E-waste Monitor, 2024)}
  \caption{Growth of Electronic Waste Management. Source: Adapted from the Global E-waste Monitor (2024).}
  \label{fig:ewaste-growth}
\end{figure}

Globally, electronic waste has emerged as a major environmental, economic, and
public health concern. Unlike ordinary municipal waste, electronic waste contains both
valuable and hazardous materials. Components such as gold, silver, copper,
aluminium, and rare earth metals can be recovered and reused through appropriate
recycling processes. However, electronic waste also contains hazardous substances,
including lead, mercury, cadmium, chromium, brominated flame retardants, and other
toxic chemicals capable of contaminating soil, water resources, and the atmosphere
when improperly handled or disposed of. Consequently, ineffective electronic waste
management contributes to environmental pollution, climate change, biodiversity
loss, and serious health risks for individuals involved in informal recycling activities.

The concept of the circular economy has gained increasing international attention
as a sustainable approach to addressing the growing electronic waste challenge.
Unlike the traditional linear economy, which follows the sequence of production,
consumption, and disposal, the circular economy promotes resource efficiency by
extending product life cycles through repair, refurbishment, reuse, remanufacturing,
and recycling. This approach minimizes waste generation while maximizing resource
recovery and environmental sustainability. International organizations such as the
United Nations, the International Telecommunication Union (ITU), and the Global
E-waste Monitor continue to advocate for technology-driven solutions that enhance
traceability, accountability, and efficient recovery of electronic waste.
Figure~\ref{fig:ewaste-lifecycle} illustrates the complete electronic waste lifecycle,
beginning with product design and use, followed by collection, transportation,
treatment, recycling, and resource recovery within a circular economy framework.

\begin{figure}[H]
  \centering
  \imageplaceholder{Global Electronic Waste Lifecycle diagram (adapted from the Global E-waste Monitor, 2024)}
  \caption{Global Electronic Waste Lifecycle. Source: Adapted from the Global E-waste Monitor (2024).}
  \label{fig:ewaste-lifecycle}
\end{figure}

Recent advances in digital technologies have transformed waste management
practices worldwide. Cloud computing, mobile applications, artificial intelligence,
geographic information systems (GIS), Internet of Things (IoT) technologies, and big
data analytics are increasingly being integrated into modern environmental
management systems. Mobile applications enable citizens to report waste locations,
schedule collection services, and receive notifications, while cloud platforms
facilitate centralized data management, scalability, and real-time information sharing.
Artificial intelligence has also demonstrated significant potential in automating waste
classification through image recognition, thereby improving sorting accuracy and
reducing manual effort. Geographic information systems further support route
planning, location tracking, and optimized collection logistics.

Despite these technological advancements, electronic waste management in many
developing countries remains largely manual and fragmented. Collection processes
often rely on informal practices with limited record-keeping, inefficient transportation
planning, inadequate monitoring mechanisms, and poor coordination among
stakeholders. Consequently, electronic waste frequently ends up in open dumpsites,
landfills, or informal recycling operations where unsafe handling practices expose
workers and surrounding communities to hazardous materials.

In Sierra Leone, the increasing adoption of information and communication
technologies has significantly increased the number of electronic devices used by
households, businesses, educational institutions, healthcare facilities, and government
organizations. However, the country's electronic waste management infrastructure
has not expanded at the same pace. Most electronic waste management activities
continue to rely on manual documentation, informal collection, limited recycling
facilities, and inadequate mechanisms for tracking electronic waste throughout its
lifecycle. Furthermore, there is limited public awareness regarding responsible
electronic waste disposal, resulting in environmentally harmful disposal practices.

The absence of integrated digital platforms further complicates electronic waste
management. Existing approaches generally lack centralized information systems
capable of registering electronic waste, tracking collection activities, coordinating
transportation, facilitating repair and refurbishment, supporting recycling operations,
or generating comprehensive environmental reports. Similarly, many existing systems
do not leverage artificial intelligence to classify electronic waste automatically or
utilize geospatial technologies to optimize collection routes and improve operational
efficiency.

Advances in cross-platform software development now provide opportunities to
address these limitations through integrated digital solutions. Frameworks such as
Flutter enable developers to build applications that operate consistently across
Android, web, and Windows platforms from a single codebase. Backend frameworks
such as FastAPI provide high-performance web services for secure communication
between client applications and cloud infrastructure. Cloud-based technologies
including Firebase Authentication, Cloud Firestore, Cloudinary, Firebase Cloud
Messaging, and Google Maps Platform further enable secure authentication, scalable
data management, cloud image storage, push notifications, and location-based
services.

Motivated by these challenges and opportunities, this project proposes the
development of EcoTrace: Intelligent Electronic Waste Recovery and Circular
Economy Management System. EcoTrace is designed as a cross-platform intelligent
information system that supports the entire electronic waste management lifecycle.
The system enables users to register electronic waste, upload device images, request
waste collection, monitor transportation activities, support repair, refurbishment, and
recycling processes, generate administrative reports, and improve stakeholder
collaboration through a centralized digital platform. The system also integrates cloud
technologies, geographic information systems, and artificial intelligence to enhance
traceability, operational efficiency, and environmental sustainability.

The development of EcoTrace aligns with global sustainable development
initiatives by promoting responsible consumption, environmental protection, and
circular economy practices through the application of modern software engineering
principles and intelligent computing technologies. The project therefore contributes
not only to improving electronic waste management processes but also to
demonstrating the practical application of cloud computing, artificial intelligence, and
cross-platform software development in addressing real-world environmental
challenges.

\section{Statement of the Problem}
\label{sec:problem-statement}

The rapid increase in the consumption of electrical and electronic equipment has
resulted in a corresponding increase in the generation of electronic waste worldwide.
While many developed countries have adopted integrated digital systems and
regulatory frameworks to support electronic waste collection, tracking, recycling, and
environmental compliance, many developing countries continue to face significant
challenges in managing electronic waste effectively. The absence of efficient
information management systems contributes to poor monitoring, limited
accountability, environmental pollution, and the loss of valuable recyclable materials.

In Sierra Leone, electronic waste management is still characterized by
fragmented and largely manual processes. Households, businesses, educational
institutions, healthcare facilities, and government agencies frequently dispose of
obsolete electronic devices without structured collection mechanisms or proper
documentation. In many cases, electronic waste is mixed with municipal solid waste,
openly burned, or handled by informal recyclers using unsafe techniques that expose
both people and the environment to hazardous substances. These practices increase
the risk of soil contamination, water pollution, air pollution, and adverse public health
effects.

Existing electronic waste management practices also suffer from limited
coordination among key stakeholders. Waste generators, collectors, transporters,
technicians, recyclers, environmental officers, and regulatory authorities often operate
independently with little or no centralized information sharing. Consequently, it
becomes difficult to monitor electronic waste from the point of generation to its final
recovery or disposal. The lack of traceability reduces transparency, limits
accountability, and makes it difficult for authorities to evaluate the effectiveness of
electronic waste management programmes.

Furthermore, the absence of intelligent digital tools limits operational efficiency.
Most existing practices do not provide automated electronic waste classification,
optimized collection routes, centralized cloud storage, real-time notifications, or
comprehensive reporting capabilities. Decision-makers therefore lack timely and
accurate information required for planning, resource allocation, policy formulation,
and environmental monitoring. Similarly, citizens have limited access to convenient
digital platforms through which they can register electronic waste, schedule
collections, or track the progress of their requests.

Although several waste management applications have been developed in
different parts of the world, many of them focus primarily on general solid waste
management rather than the complete lifecycle of electronic waste. In addition, many
existing systems do not integrate artificial intelligence for automated electronic waste
classification, cloud-based cross-platform accessibility, image management,
geospatial services, and end-to-end lifecycle tracking within a single platform. These
limitations reduce their suitability for addressing the specific electronic waste
management challenges experienced in Sierra Leone.

The absence of a comprehensive, intelligent, and integrated electronic waste
management platform therefore creates significant challenges in electronic waste
registration, collection, transportation, repair, refurbishment, recycling, monitoring,
reporting, and environmental decision-making. There is consequently a need for a
modern digital solution capable of supporting the complete electronic waste
management lifecycle while improving operational efficiency, stakeholder
collaboration, transparency, traceability, and environmental sustainability.

This study addresses these challenges through the design and implementation of
the EcoTrace: Intelligent Electronic Waste Recovery and Circular Economy
Management System, which integrates cross-platform mobile, web, and desktop
applications with cloud computing, artificial intelligence, geospatial technologies, and
secure backend services to improve the management of electronic waste in Sierra
Leone.

\section{Aim and Objectives of the Study}
\label{sec:aim-objectives}

\subsection{Aim of the Study}

The main aim of this study is to design, develop, and implement an intelligent
cross-platform electronic waste recovery and circular economy management system
that leverages cloud computing, artificial intelligence, and geospatial technologies to
improve the registration, collection, tracking, recovery, and sustainable management
of electronic waste in Sierra Leone.

\subsection{Specific Objectives}

The specific objectives of this study are to:

\begin{enumerate}[label=\arabic*.,leftmargin=1.4cm]
  \item Design and develop a secure cross-platform electronic waste management
    system that operates on Android, Web, and Windows platforms.
  \item Develop a centralized cloud-based platform for registering, storing, managing,
    and tracking electronic waste information using Cloud Firestore.
  \item Implement secure user authentication and role-based access control for
    different categories of system users using Firebase Authentication.
  \item Integrate artificial intelligence to support the automatic classification of
    electronic waste from uploaded images.
  \item Develop an intelligent electronic waste pickup scheduling and collection
    management module for households and business organizations.
  \item Integrate Google Maps Platform to support location services, display
    collection locations, visualize routes, and improve collection operations.
  \item Develop modules for electronic waste transportation, repair, refurbishment,
    recycling, and lifecycle tracking to promote circular economy practices.
  \item Implement real-time notification services using Firebase Cloud Messaging to
    keep users informed about collection requests, status updates, and system
    activities.
  \item Generate administrative reports and analytics to support environmental
    monitoring, operational planning, and informed decision-making.
  \item Evaluate the effectiveness, usability, and performance of the developed
    EcoTrace system in improving electronic waste management processes.
\end{enumerate}

\section{Significance of the Study}
\label{sec:significance}

This study is significant because it provides a practical technological solution for
improving electronic waste management through the integration of cloud computing,
artificial intelligence, geospatial technologies, and cross-platform software
development. The developed EcoTrace system is expected to benefit various
stakeholders involved in electronic waste generation, collection, recovery,
environmental management, policy formulation, research, and sustainable
development.

\subsection{Households and Individual Users}
EcoTrace provides households and individual users with a convenient and secure
platform for registering electronic waste, scheduling collection requests, monitoring
collection progress, and receiving real-time notifications. The system encourages
responsible disposal practices and improves public participation in environmental
conservation.

\subsection{Business Organizations}
Business organizations frequently generate large quantities of obsolete electronic
equipment. EcoTrace enables organizations to maintain proper records of electronic
waste, schedule collections efficiently, monitor asset disposal, and ensure
environmentally responsible management of obsolete electronic devices.

\subsection{Electronic Waste Collectors}
The system assists electronic waste collectors by providing organized pickup
requests, optimized collection schedules, customer location information, and
real-time status updates. These features improve operational efficiency while
reducing unnecessary travel time and fuel consumption.

\subsection{Drivers and Logistics Personnel}
Through the integration of Google Maps Platform, drivers can visualize collection
locations, follow optimized routes, and update the status of transportation activities.
This contributes to more efficient logistics, reduced operational costs, and improved
coordination during electronic waste collection.

\subsection{Technicians and Refurbishment Centres}
Technicians responsible for inspecting, repairing, and refurbishing electronic devices
benefit from centralized records that support the tracking of repair activities,
refurbishment status, and equipment history. This promotes accountability while
extending the useful life of electronic devices in accordance with circular economy
principles.

\subsection{Recycling Companies}
Recycling organizations benefit from improved traceability of electronic waste
received for processing. The system maintains digital records of recycling activities,
enabling better inventory management, regulatory reporting, and resource recovery.

\subsection{Environmental Protection Agencies}
Environmental officers and regulatory institutions can use EcoTrace to monitor
electronic waste generation, collection, transportation, recycling activities, and
environmental performance. Administrative reports and analytical dashboards support
evidence-based decision-making, policy implementation, and environmental
compliance monitoring.

\subsection{Government and Policy Makers}
Government institutions can utilize the information generated by the system to
formulate policies, evaluate electronic waste management programmes, allocate
resources effectively, and develop national strategies that promote environmental
sustainability and responsible electronic waste management.

\subsection{Researchers and Academic Institutions}
The study contributes to existing knowledge in software engineering, environmental
informatics, and intelligent waste management by demonstrating the practical
integration of cloud computing, artificial intelligence, mobile computing, geospatial
technologies, and modern software development frameworks. The project may also
serve as a useful reference for future academic research in related fields.

\subsection{The Computer Science Profession}
The project demonstrates the application of software engineering principles to
solving a real-world environmental problem. It highlights how modern computing
technologies can be integrated to develop scalable, secure, and intelligent
information systems capable of supporting sustainable development objectives.

\subsection{The General Public}
The successful implementation of EcoTrace is expected to improve public awareness
of responsible electronic waste disposal, reduce environmental pollution, encourage
recycling and resource recovery, and contribute to the protection of public health. By
promoting environmentally responsible practices, the system supports sustainable
development and the long-term conservation of natural resources.

\section{Scope of the Study}
\label{sec:scope}

The scope of this study covers the analysis, design, development, implementation,
and evaluation of the EcoTrace: Intelligent Electronic Waste Recovery and
Circular Economy Management System. The study focuses on developing an
integrated cross-platform solution that supports the efficient management of
electronic waste throughout its lifecycle, from registration and collection to
transportation, repair, refurbishment, recycling, and final reporting.

The developed system is implemented as a cross-platform application using
Flutter, enabling deployment on Android mobile devices, web browsers, and Windows
desktop computers from a unified codebase. The backend services are implemented
using FastAPI and deployed on the Render cloud platform to provide secure, scalable,
and high-performance application programming interfaces (APIs). Cloud Firestore
serves as the primary cloud database for storing system data, while Firebase
Authentication provides secure user authentication and role-based access control.
Cloudinary is used for cloud-based image storage, Firebase Cloud Messaging delivers
push notifications, and Google Maps Platform supports location services, route
visualization, and collection route optimization. Artificial intelligence functionality is
implemented using PyTorch to support electronic waste image classification.
Chapter~\ref{ch:implementation} documents the technology stack exactly as
implemented, including the points at which the delivered system took a different
technical path from what is described here.

The system supports multiple categories of users, including household users,
business organizations, electronic waste collectors, drivers, technicians, recyclers,
environmental officers, administrators, and super administrators. Each user category
is provided with specific functionalities based on assigned roles and responsibilities
within the electronic waste management process.

The functional scope of EcoTrace includes user registration and authentication,
electronic waste registration, image upload, artificial-intelligence-based waste
classification, pickup request management, collection scheduling, transportation
monitoring, repair and refurbishment management, recycling management, location
visualization, notification services, reporting, analytics, and administrative
management.

Geographically, the study is designed with Sierra Leone as its primary
application context. However, the architecture and technologies used are designed to
be scalable and adaptable for deployment in other regions with similar electronic
waste management requirements.

This study does not include the physical construction or operation of electronic
waste recycling facilities, the manufacturing of recycling equipment, or the direct
processing of hazardous electronic waste. Similarly, the study does not evaluate
national environmental policies or perform large-scale environmental impact
assessments. Instead, its primary focus is the development and evaluation of an
intelligent software platform that supports efficient electronic waste management
through modern information and communication technologies.

\section{Limitations of the Study}
\label{sec:limitations-intro}

Although the EcoTrace system was successfully designed and implemented to
support intelligent electronic waste management, the study was subject to several
limitations that may influence the deployment, evaluation, and future enhancement of
the system.

Firstly, the backend services were deployed using the free tier of the Render
cloud platform. While this deployment environment is adequate for development,
demonstration, and small-scale usage, it provides limited computing resources and
may experience delayed response times due to automatic service suspension during
periods of inactivity. Consequently, the system may not provide the same level of
performance expected from a production-grade deployment hosted on dedicated
cloud infrastructure.

Secondly, the operation of the system depends on a stable Internet connection
because cloud-based services are used for authentication, database management,
image hosting, notifications, and backend communication. Users operating in areas
with unreliable or limited Internet connectivity may experience delays in data
synchronization, image uploads, and real-time system updates.

Another limitation relates to the artificial intelligence component. Although the
system integrates PyTorch to support intelligent classification of electronic waste, the
performance of the AI model is dependent on the quality, diversity, and quantity of the
training dataset. Variations in lighting conditions, camera quality, image resolution,
and device orientation may affect classification accuracy, particularly for damaged,
uncommon, or visually similar electronic devices.

The system also relies on several third-party cloud services, including Cloud
Firestore, Firebase Authentication, Firebase Cloud Messaging, Cloudinary, Google
Maps Platform, and Render. Any temporary service interruptions, API changes,
pricing adjustments, or usage limitations introduced by these service providers may
affect the availability or performance of specific system functionalities.

Furthermore, this study focuses primarily on the software engineering aspects of
electronic waste management. The project does not include the construction or
operation of physical recycling facilities, laboratory testing of hazardous electronic
materials, or large-scale environmental impact assessments. Consequently, the
evaluation of the system is based on software functionality, usability, and system
performance rather than the direct measurement of environmental outcomes.

Finally, the study was developed using Sierra Leone as the primary case study.
Although the overall system architecture is scalable and can be adapted for
deployment in other countries, certain operational procedures, institutional
arrangements, and regulatory requirements may vary across jurisdictions and would
require appropriate customization before large-scale implementation.

\section{Definition of Terms}
\label{sec:definitions}

For the purpose of this study, the following terms are defined as follows:

\begin{description}[leftmargin=0cm,style=nextline]
  \item[Electronic Waste (E-waste)] Discarded electrical and electronic equipment
    that has reached the end of its useful life or is no longer required by its owner.
  \item[Circular Economy] An economic model that promotes the reuse, repair,
    refurbishment, remanufacturing, and recycling of products and materials to
    minimize waste generation and maximize resource utilization.
  \item[Waste Electrical and Electronic Equipment (WEEE)] Electrical and
    electronic devices that have become obsolete, damaged, or unwanted and require
    proper collection, treatment, recycling, or disposal.
  \item[Artificial Intelligence (AI)] A branch of computer science that enables
    computer systems to perform tasks that normally require human intelligence,
    such as learning, classification, prediction, and decision-making.
  \item[Machine Learning] A subset of artificial intelligence that enables computer
    systems to learn patterns from data and improve performance without being
    explicitly programmed for every task.
  \item[PyTorch] An open-source machine learning framework used in this study to
    develop and deploy the electronic waste image classification model.
  \item[Cloud Computing] The delivery of computing services such as servers,
    databases, networking, storage, and software over the Internet.
  \item[Cloud Firestore] A cloud-hosted NoSQL database provided by Firebase for
    storing, synchronizing, and managing application data in real time.
  \item[Firebase Authentication] A cloud-based authentication service used to
    securely identify users and manage access to system resources.
  \item[Firebase Cloud Messaging (FCM)] A cloud messaging service that enables
    the delivery of real-time push notifications to registered application users.
  \item[Cloudinary] A cloud-based media management platform used for storing,
    managing, and delivering electronic waste images uploaded by users.
  \item[FastAPI] A modern Python web framework used to develop high-performance
    backend Application Programming Interfaces (APIs) for the EcoTrace system.
  \item[Flutter] Google's open-source cross-platform software development
    framework used to develop Android, Web, and Windows applications from a
    single codebase.
  \item[Google Maps Platform] A collection of geospatial services used for
    displaying maps, visualizing collection locations, and supporting route planning
    within the EcoTrace system.
  \item[Route Optimization] The process of determining the most efficient travel
    route between multiple collection locations to reduce travel distance, time, and
    operational cost.
  \item[Application Programming Interface (API)] A set of rules and protocols
    that enables communication between software applications and backend
    services.
  \item[Role-Based Access Control (RBAC)] A security mechanism that restricts
    system access based on the roles and responsibilities assigned to authenticated
    users.
  \item[Cross-Platform Application] A software application developed once and
    deployed across multiple operating systems with minimal platform-specific
    modifications.
\end{description}

\section{Organization of the Report}
\label{sec:organization}

This report is organized into five chapters, each addressing a specific aspect of the
study.

Chapter One introduces the study by presenting the background of the study,
statement of the problem, aim and objectives, significance of the study, scope,
limitations, definition of terms, and the organization of the report.

Chapter Two reviews relevant literature related to electronic waste management,
circular economy principles, cloud computing, artificial intelligence, cross-platform
application development, and existing electronic waste management systems. It also
identifies the research gap addressed by this study.

Chapter Three describes the research methodology and presents the analysis
and design of the proposed EcoTrace system. It discusses the existing system, system
requirements, software architecture, database design, system models, user interface
design, and the technologies used in the development process.

Chapter Four presents the implementation and testing of the EcoTrace system.
It describes the development environment, system modules, implementation
procedures, testing activities, evaluation results, and screenshots of the developed
application.

Chapter Five summarizes the study, presents the conclusions drawn from the
research findings, highlights the contribution of the study, provides recommendations
for stakeholders, and suggests areas for future research and system enhancement.

\section{Chapter Summary}

This chapter introduced the study by discussing the background of electronic waste
management, identifying the research problem, presenting the aim and objectives of
the study, describing its significance, defining its scope and limitations, explaining
key terms, and outlining the organization of the report. The next chapter reviews
relevant literature, existing technologies, and related electronic waste management
systems that provide the theoretical and technological foundation for the
development of EcoTrace.
